\chapter{Az adatbázis-engineering kiegészítő alprojekt: databaseproject és C\# böngésző}

\section{A kiegészítő alprojekt szerepe és motivációja}

A Premade Graph fő rendszere SQLite-alapú perzisztenciával dolgozik, és a gráfanalitikai pipeline-hoz és a frontendhez optimalizált sémát használ. A \texttt{docs/databaseproject} könyvtárban egy ezzel párhuzamos, de önálló alprojekt él: egy formálisan megtervezett, PostgreSQL-re épülő relációs adatbázismodell és egy ASP.NET Core / C\# alapú adatbázis-böngésző alkalmazás. Ez az alprojekt kettős értékkel bír: egyrészt erősíti az adatbázis-tervezés akadémiai oldalát, másrészt valódi, futtatható eszközt biztosít az elemzési eredmények böngészéséhez.

Az \texttt{architecture-decision} szemlélet alapján fontos tisztázni, miért kettő rendszer létezik egymás mellett. Az SQLite-alapú fő rendszer a gráfépítés, az analitika és a demo-pipeline kiszolgálásához optimalizált: egyszerű telepítés, nincs szerverfüggőség, gyors az iteráció. A PostgreSQL-alapú adatbázismodell ezzel szemben a relációs adatbázistervezés akadémiai igényeivel igazodik: tárolt eljárások, triggerek, rekurzív lekérdezések, JSON-mezők, hivatkozási integritás. A kettő nem versenyez egymással, hanem különböző célt szolgál ugyanabban a domainben.

\section{A formális relációs séma}

A \texttt{docs/databaseproject/schema.sql} egy teljes PostgreSQL sémát definiál. A séma tíz főtáblája a következő:

\begin{table}[H]
\centering
\caption{A databaseproject PostgreSQL sémájának táblái}
\begin{tabularx}{\textwidth}{>{\raggedright\arraybackslash}p{4.5cm}X}
\toprule
\textbf{Tábla} & \textbf{Szerepkör a sémában} \\
\midrule
\texttt{players} & Játékosazonosság, összesített pontszámok, szerepkör \\
\texttt{matches} & Meccsszintű adatok, queue, patch, időtartam, győztes \\
\texttt{teams} & Egy meccs két csapata, oldal, objektíva adatok \\
\texttt{match\_participants} & Résztvevői adatok: champion, lane, kills, deaths, assists \\
\texttt{player\_relationships} & Játékos--játékos ally/enemy kapcsolat súllyal \\
\texttt{clusters} & Klaszter metaadatok, típus, algoritmus \\
\texttt{cluster\_members} & Tagsági rekordok, szerepkör-összegzés \\
\texttt{analysis\_runs} & Elemzési futások metaadatai \\
\texttt{performance\_snapshots} & Játékos-teljesítmény pillanatfelvételek futásonként \\
\texttt{path\_queries} & Útkeresési kérések és eredmények \\
\bottomrule
\end{tabularx}
\end{table}

A séma normalizált: a meccsadatok nem ismétlődnek játékosonként, a résztvevői rekord tartalmazza a meccsre vonatkozó joinolt mezőket, a klasztertagság idegen kulcson keresztül kapcsolódik a játékoshoz és a klaszterhez. A \texttt{CHECK} constraintek biztosítják, hogy az \texttt{opscore} és \texttt{feedscore} értékek a $[0,10]$ tartományban maradnak.

\subsection{Nézetek és összetett lekérdezések}

A \texttt{docs/databaseproject/views\_queries.sql} fájl elemző nézeteket és lekérdezéseket tartalmaz. Ezek közé tartoznak:
\begin{itemize}
    \item \textbf{Szerepkörös teljesítményvizsgálat}: aggregált opscore/feedscore elemzés szerepkörönként;
    \item \textbf{Klaszterösszegzők}: belső sűrűség, hídjátékos-jelzők, átlagos teljesítmény klaszterenként;
    \item \textbf{Kapcsolat-erősségi rang}: a legerősebb ismétlődő ally/enemy pár keresése;
    \item \textbf{Teljesítménystabilitás időben}: teljesítmény-pillanatfelvételek összehasonlítása elemzési futásonként;
    \item \textbf{JSON-alapú metaadat-lekérdezések}: a klaszter \texttt{summary\_json} és \texttt{role\_json} mezőinek lekérdezése PostgreSQL JSON operátorokkal.
\end{itemize}

\subsection{Tárolt eljárások és triggerek}

A \texttt{docs/databaseproject/procedures\_triggers.sql} tárolt eljárásokat és triggereket definiál. Ezek közé tartoznak:
\begin{itemize}
    \item \textbf{Pontszámfrissítési trigger}: meccs-résztvevő beírásakor automatikusan frissíti az aggregált player-szintű opscore-t;
    \item \textbf{Klaszterépítési eljárás}: adott elemzési futáshoz klaszterező logikát futtat és klasztertagságot rögzít;
    \item \textbf{Útkeresés-naplózó eljárás}: útkeresési kérés és eredmény konzisztens rögzítése;
    \item \textbf{Teljesítménypillanatkép-karbantartó}: elemzési futásonként snapshot-ot ír, megőrizve a longitudinális adatot.
\end{itemize}

Az akadémiai szempont itt az, hogy a tárolt eljárások és triggerek az adatbázis-logikát az alkalmazáslogikától elválasztják: az adatmodellek önállóan karbantarthatók, és az elemzési pipeline nem függhet attól, hogy az alkalmazásréteg helyesen hívja-e meg az aggregálási logikát.

\section{Az ASP.NET Core / C\# adatbázis-böngésző}

A \texttt{docs/databaseproject/simple-interface/} könyvtárban egy ASP.NET Core webalkalmazás él, amelyet a Premade Graph SQLite adatbázisainak böngészésére terveztek. A böngésző \texttt{Program.cs} fájlja mutatja, hogyan épül fel ez az eszköz.

\subsection{Adatbázis-generálás és importálás}

Az alkalmazás indulásakor automatikusan futtatja a \texttt{ProjectImporter.Import()} folyamatot, amely a Premade Graph repositoryban lévő SQLite adatbázisból importál a böngésző saját, projektre generált SQLite állományába:

\begin{itemize}
    \item \texttt{players}: opscore, feedscore, szerepkör, meccsszám;
    \item \texttt{clusters}: klasztertípus, méret, metaadatok;
    \item \texttt{cluster\_members}: tagsági rekordok;
    \item \texttt{replay}: útkeresési mentések;
    \item \texttt{replay\_run}: útkeresési futások.
\end{itemize}

Az import eredménye visszajelzést ad a konzolra és webes hibaoldalon jelzi, ha az importált adatbázis nem elérhető. Ez azt jelenti, hogy a böngésző önállóan indítható, és ha az adatforrás hiányzik, explicit hibaüzenetet ad, nem csendesen üres felületet mutat.

\subsection{Dataset-szelektálás}

Az alkalmazás tudatosan illeszkedik a Premade Graph több-adatkészletes architektúrájához. A \texttt{DatasetCatalog.Load(root)} a repository \texttt{datasets.json} regiszterét olvassa be, és a böngésző minden oldalán URL-paraméterként átadható a kívánt dataset azonosítója (\texttt{?dataset=flexset}, \texttt{?dataset=soloq}).

Ez azt jelenti, hogy egyetlen böngészőpéldány az összes regisztrált adatkészlethez adatbázist tud generálni és megjeleníteni, anélkül hogy a mögöttes SQLite állományokat összekeverne.

\subsection{Táblabögészés, szűrés, lapozás és JSON-nézetek}

A böngésző a következő funkciókat valósítja meg:

\begin{itemize}
    \item \textbf{Táblalistázás}: az összes importált tábla nevének és sorainak megjelenítése;
    \item \textbf{Szűrés}: szövegmező alapú szűrés játékosnévre, szereplőre, klasztertípusra;
    \item \textbf{Lapozás}: nagy táblák esetén oldalanként meghatározott számú sor;
    \item \textbf{Rendezés}: oszlop szerinti növekvő/csökkenő sorrend;
    \item \textbf{Teljes értéknézet}: a \texttt{/database/value} endpoint egy adott cella teljes értékét mutatja meg, különösen hasznos a \texttt{summary\_json} és \texttt{role\_json} típusú, hosszú JSON-oszlopoknál.
\end{itemize}

\subsection{Technikai architektúra}

Az alkalmazás egész programja egyetlen \texttt{Program.cs} fájlban helyezkedik el, minimális top-level statement stílussal. Az ASP.NET Core pipeline konfigurálása, az útvonaldefiníciók és az adatbázis-lekérdezések egységes szerkezetben élnek. A renderelés szerver oldali HTML-generálással történik, nincs JavaScript frontend: a böngésző célja az áttekinthetőség és a bemutathatóság, nem egy SPA-szerű élmény.

A \texttt{Microsoft.Data.Sqlite} NuGet csomag az adatbázis-elérés alapja, amely a SQLite adatbázisokhoz natív, közvetlen hozzáférést biztosít .NET-ből. A kapcsolatfelépítés és a lekérdezések a \texttt{ProjectDatabase} osztályon keresztül vannak absztraháLva, ami megkönnyíti a dataset-váltás és az import-logika szétválasztását.

\section{A kiegészítő alprojekt helye a dolgozatban}

\subsection{Miért valódi alprojekt, nem csak dokumentáció?}

A databaseproject nem puszta papír-terv: a \texttt{schema.sql} futtatható, a \texttt{views\_queries.sql} lekérdezései tesztelhetők PostgreSQL-ben, a C\# böngésző lefordul és futtatható a repository-ból. Mindez valódi, védhető mellékprojektté teszi.

Az akadémiai értéke kettős:
\begin{itemize}
    \item \textbf{Adatbázis-tervezési szempontból}: bemutatja, hogyan formalizálható ugyanaz a domainmodell egy erősebb, tárolt logikát támogató adatbázisrendszerben; ez az SQLite-on dolgozó fő rendszer korlátait is megvilágítja.
    \item \textbf{Szoftverarchitektúra szempontból}: demonstrálja, hogy a Premade Graph domain-modellje természetesen leírható relációs sémával, tárolt eljárásokkal és triggerekkel, miközben az alapvető fogalmak (players, clusters, cluster\_members, opscore, feedscore) mindkét adatbázis-környezetben koherensek.
\end{itemize}

\subsection{A C\# böngésző mint prezentációs eszköz}

A C\# böngésző azért is értékes a dolgozat szempontjából, mert lehetővé teszi, hogy a SQLite-ban tárolt elemzési eredmények közvetlenül, táblázatosan és JSON-nézettel ellenőrizhetők legyenek. Ez a reprodukálható kutatás egyik alapkövetelménye: az eredmények nemcsak a frontend vizualizációján láthatók, hanem adatbázis-szinten is auditálhatók \cite{sandve2013}.

\subsection{A databaseproject korlátai}

A databaseproject PostgreSQL-sémája és a fő rendszer SQLite-adatbázisa között nincs automatikus, élő szinkronizáció. A PostgreSQL-modell akadémiai-demonstrációs célt szolgál, nem production-szinkronizációt. Ez explicit módon van dokumentálva a \texttt{docs/databaseproject/documentation.tex} fájlban: ``A beadandó adatbázis-változat PostgreSQL-re készült, mert a feladat tárolt eljárásokat, triggereket, rekurzív lekérdezéseket, JSON mezőket és összetettebb SQL elemzéseket is kér. A tényleges alkalmazás jelenleg SQLite-ot használ helyi perzisztenciára, de a dokumentált modell ugyanazokat a fogalmakat viszi át egy erősebb adatbázisrendszerbe.''

Ez a szétválasztás előny, nem hiba: megmutatja, hogy a szoftvermérnöki döntés (SQLite a fő rendszerben) indokolt a project méretéhez és céljához, miközben az akadémiai adatbázistervezési kompetencia bemutatható a PostgreSQL-modellen keresztül.
